\chapter{Polynômes, racines et factorisation dans le plan complexe}
\section*{Exercice 27.1 -- Second degré réel}
Le discriminant est $4-20=-16$. Les racines sont $(2\pm4i)/2=1\pm2i$. On contrôle leur somme $2$ et leur produit $5$.

\section*{Exercice 27.2 -- Factorisation}
$P(z)=(z-1)(z^2+z+1)$. Les racines sont $1$ et $(-1\pm i\sqrt3)/2$. Elles sont distinctes et forment les trois racines cubiques de l'unité.

\section*{Exercice 27.3 -- Racine double}
La division par $z-1$ donne $z^2+z-2=(z-1)(z+2)$. Donc $P(z)=(z-1)^2(z+2)$ : deux racines distinctes, $1$ de multiplicité $2$ et $-2$ simple. Le degré compte les multiplicités dans une factorisation complète.

\section*{Exercice 27.4 -- Racines conjuguées}
La conjugaison préserve sommes et produits et fixe les coefficients réels. Donc $P(\overline z)=\overline{P(z)}$. Si $P(z)=0$, le membre de droite est nul. La condition sur les coefficients est essentielle : $P(z)=z-i$ ne possède pas $-i$ pour racine.

\section*{Exercice 27.5 -- Nombre de racines}
Un polynôme non nul de degré $d$ possède au plus $d$ racines distinctes : chaque racine se factorise et diminue le degré d'une unité. Si le polynôme proposé était non nul, ses cinq racines contrediraient $d\le4$.

\section*{Exercice 27.6 -- Racine complexe d’un complexe}
L'identification donne $x^2-y^2=3$, $2xy=4$. De plus $x^2+y^2=|z|^2=|3+4i|=5$. Ainsi $x^2=4$, $y^2=1$ et $xy=2>0$. Les solutions sont $2+i$ et $-2-i$.
